\section{Heterogeneity and distributional effects}\label{sec:heterogeneity}

\subsection{Who adjusts?}

Figure~\ref{fig:het} reports the low- versus high-skilled DiD estimate for each headline
outcome within demographic and regional subgroups.\footnote{Each estimate is the
within-subgroup difference-in-differences, so it asks whether the low- versus high-skilled
contrast differs across, say, men and women, not the level effect for a subgroup.} Three
patterns stand out and are
consistent with economic intuition about where the minimum wage binds hardest.

First, no subgroup shows wage gains. Whether defined by sex, age, area, or regional bite,
every subgroup wage estimate is negative or indistinguishable from zero. The absence of a
pay response is a pervasive feature of the data rather than an average that masks
offsetting effects.

Second, the compositional adjustment scales with exposure. The decline in formal
employment is larger in high-bite departments ($-5.5$ percentage points) than in low-bite
departments ($-3.6$ points), and the rise in self-employment likewise ($+5.0$ against
$+2.8$ points). This dose-response pattern is what a causal effect of the wage floor
should produce, and it is the cross-sectional counterpart of the continuous-exposure
results in Section~\ref{sec:robustness}. The employment-level estimate, by contrast, does
\emph{not} scale with the bite---it is smaller and insignificant in high-bite
departments---which is further reason to attribute it to differential recovery rather
than to the reform.

Third, the reallocation is concentrated among the young and among women. The decline in
formal employment is largest for workers aged fourteen to twenty-five ($-5.8$ points) and
for women ($-5.5$ points), the groups whose wages are most likely to sit near the floor
and whose attachment to formal employment is most tenuous; the rise in self-employment is
twice as large for women as for men. For workers over forty-five the formality effect is
near zero, though their self-employment also rises. This mirrors the emphasis on youth in
the earlier Peruvian literature \citep{cespedes_2005, jaramillo_lopez_2006} and the
emphasis on women and the least skilled in the developing-country evidence
\citep{katzkowicz_2021_uruguay, lemos_2009}.

\begin{figure}[!t]
\centering
\includegraphics[width=0.98\textwidth]{fig6_heterogeneity.pdf}
\caption{Heterogeneous effects of the 2022 reform. Each point is the
$\mathrm{Low}\times\mathrm{Post}$ DiD estimate within the indicated subgroup, with
ninety-five percent confidence intervals based on standard errors clustered by department.}
\label{fig:het}
\end{figure}

\subsection{Adjustment margins}

Where exactly does the reallocation happen? Four mechanism-oriented margins sharpen
the picture (Appendix Table~\ref{tab:margins}). First, the share of low-education
wage earners employed \emph{without a written contract} rises by 4.5 percentage
points relative to high-education wage earners---informalization operates within
wage employment, not only through exits into self-employment. Second, employment
recomposes toward \emph{micro units}: the probability that a low-education worker is
employed in a establishment of at most twenty workers rises by 2.8 points, and small
units are precisely where enforcement is thinnest. Third, using INEI's decomposition
of informal employment (available through 2023), the displaced margin lands in the
\emph{informal sector} proper (+2.4 points) rather than in unregistered jobs inside
formal firms (+0.8, insignificant): workers did not stay with formal employers under
the table; they moved to informal productive units. Fourth, the share of wage
earners with under a year of tenure does not move, so the adjustment is not visible
as a hiring freeze among those who remain wage earners.

Two further subgroup splits locate the burden. The compositional effects are
concentrated almost entirely among \emph{non-heads} of household---secondary
earners, whose formality falls 6.8 points against a null for household heads---and
among non-indigenous workers, consistent with the urban coastal labor markets where
formal employment at the floor is concentrated. The marginal formal jobs that
dissolved were, in short, those of secondary earners in small urban firms.

\subsection{Distributional wage effects}

The absence of an average wage effect could in principle conceal offsetting movements at
different points of the wage distribution, since a minimum wage is expected to compress the
lower tail. To examine this I estimate unconditional quantile (recentered influence
function) regressions in the manner of \citet{firpo_fortin_lemieux_2009}, applying the DiD
specification to the influence function of each decile of the log real hourly wage.
Figure~\ref{fig:rif} plots the $\mathrm{Low}\times\mathrm{Post}$ coefficient across
deciles. The estimates are non-positive at every decile, mostly between $-2$ and $-4$
percent and significant at a few middle deciles; crucially, there is no positive effect at
the lower deciles, where a binding and enforced minimum wage would compress the
distribution from below. The distributional evidence therefore reinforces the absence of
a pay response: the reform did not lift the lower tail of the low-skilled wage
distribution relative to the high-skilled. This is consistent with the earnings
distribution in Figure~\ref{fig:dist}, which shows little movement toward the new floor,
and with the interpretation that in a labor market with widespread non-compliance the
statutory minimum did not bind for most low-skilled workers.

\begin{figure}[!t]
\centering
\includegraphics[width=0.72\textwidth]{fig8_rif_quantile.pdf}
\caption{Distributional wage effects. Each point is the $\mathrm{Low}\times\mathrm{Post}$
DiD coefficient from an unconditional-quantile (RIF) regression at the indicated decile of
the log real hourly wage \citep{firpo_fortin_lemieux_2009}. Bands are ninety-five percent
confidence intervals clustered by department.}
\label{fig:rif}
\end{figure}
